3.2 Grid Current Modeling
3.2.1 Physical Background
In a triode or pentode vacuum tube operating under normal conditions, the control grid is biased negatively with respect to the cathode, so that it draws no current. However, when the instantaneous grid-to-cathode voltage $V_{gk}$ exceeds zero — a situation that occurs during the positive half-cycles of a sufficiently large input signal — the grid begins to conduct, drawing a positive current $i_g$ through the grid resistor $R_g$. This current has two immediate consequences that profoundly shape the sonic character of the amplifier stage:

Input attenuation (blocking distortion). The voltage drop across $R_g$ partially cancels the positive excursion of the signal, effectively reducing the amplitude of the signal reaching the grid. For sustained loud playing, the coupling capacitor $C_c$ charges through $R_g$, shifting the grid bias toward negative values and causing a progressive strangling of the signal that recovers only slowly once the input level drops.

Bias point shift. The same charge accumulation shifts the quiescent operating point of the triode, altering the asymmetry of the transfer function and thereby changing the harmonic content of the output. In practice this manifests as an increase in even-order harmonics — particularly H2 — giving a dynamically varying warmth that is perceived as organic and responsive to the player's touch.

These phenomena are collectively referred to as blocking distortion in the literature [CITE], and represent one of the most distinctive and musically desirable characteristics of tube preamp stages.

3.2.2 Proposed Model
We propose a lightweight, real-time grid current model that captures both effects using only four parameters, making it suitable for interactive Web Audio applications running at standard sample rates (44.1–48 kHz).

Let $x[n]$ denote the input signal sample at discrete time $n$. The model proceeds in three stages:

Step 1 — Threshold rectification.

$$ i_g[n] = \max!\bigl(0,; x[n] - V_{th}\bigr) \tag{1} $$

where $V_{th} \in [0.05, 0.90]$ is the grid conduction threshold, normalized to the signal amplitude. A high $V_{th}$ (e.g.\ 0.7–0.8) models a clean amplifier such as the Fender Twin Reverb, where grid conduction is rarely triggered; a low $V_{th}$ (e.g.\ 0.25–0.35) models high-gain amplifiers (Peavey 5150, Mesa Dual Rectifier) where grid current is triggered frequently and aggressively.

Step 2 — RC envelope following.

The instantaneous grid current $i_g[n]$ is smoothed by a first-order IIR filter that models the $R_g C_c$ coupling network:

$$ e[n] = e[n-1] + \alpha \cdot \bigl(i_g[n] - e[n-1]\bigr), \quad \alpha = \frac{1}{\max(1,; \tau \cdot f_s \cdot 10^{-3})} \tag{2} $$

where $\tau$ (ms) is the RC time constant and $f_s$ is the sample rate. This single-pole filter efficiently approximates the charge–discharge behavior of the coupling capacitor. Short time constants ($\tau \approx 5$–20 ms) model circuits with small $R_g C_c$ products, yielding a tight, percussive response to transients; long time constants ($\tau \approx 150$–500 ms) reproduce the slow breathing associated with large coupling capacitors (e.g.\ Vox Class-A stages).

Step 3 — Dual-path signal modification.

The envelope $e[n]$ modifies the input signal through two independent paths controlled by scalar gains $\alpha_{in}$ and $\alpha_{bias}$:

$$ x_{mod}[n] = x[n] - \alpha_{in} \cdot e[n] \tag{3} $$

$$ \Delta b[n] = -\alpha_{bias} \cdot e[n] \tag{4} $$

where $x_{mod}$ is the attenuated input and $\Delta b$ is the dynamic bias shift. Both are then passed to the nonlinear triode transfer function:

$$ y[n] = \tanh!\bigl(a \cdot x_{mod}[n] + b + \Delta b[n]\bigr) \tag{5} $$

where $a$ is the tube drive gain and $b$ is the static bias offset. The complete signal processing graph is implemented in FAUST as a two-output block (inMod, biasShift) that enables efficient dataflow compilation to WebAssembly.

3.2.3 Parameter Analysis and Figure Descriptions
Figure 1 — Grid Current Characteristic. Figure 1(a) shows the rectification defined by Equation (1) for four threshold values. The piecewise-linear shape — zero for $x < V_{th}$, linearly increasing otherwise — provides a computationally efficient yet physically accurate proxy for the onset of grid conduction, which in a real tube is an exponential function of $V_{gk}$ above zero. The threshold $V_{th}$ acts as a sensitivity control: lowering it makes the stage react to a wider range of input amplitudes and triggers the effect more frequently. Figure 1(b) illustrates the input attenuation (Equation 3) for instantaneous grid current; note that the modified input $x_{mod}$ becomes increasingly nonlinear with respect to the original $x$ as $\alpha_{in}$ grows, introducing an additional soft-clipping mechanism that acts only on positive peaks.

Figure 2 --- Envelope Dynamics and Musical "Breathing". Figure 2 illustrates how the time constant $\tau$ governs the temporal "feel" of the amplifier. Panel (a) shows a rhythmic sequence of pulses (simulating palm-muted playing) that exceed the grid conduction threshold $V_{th}$. Panel (b) compares the envelope response for two distinct circuit configurations: a "tight" setting ($\tau = 15$ ms) that tracks individual hits, and a "breathing" setting ($\tau = 150$ ms) where the envelope accumulates across the sequence. Panel (c) reveals the audible impact of this accumulation through the resulting gain modulation profile: while the fast $\tau$ provides a percussive response, the slow $\tau$ creates a progressive gain reduction (analogous to power supply "sag") that experienced players describe as the amplifier "breathing" with their performance.

Figure 3 --- Grid Bias Bloom and Harmonic Evolution. Figure 3 provides a multi-stage visualization of the "harmonic bloom" effect. Panel (a) shows the internal DC bias "crawl" over 250 ms as the signal's positive peaks trigger grid conduction and charge the coupling capacitor. Panel (b) illustrates the resulting evolution of the output waveform: as the bias becomes more negative, the clipping becomes increasingly asymmetric, shifting from a relatively symmetric initial state to a "colder" clipped profile. Finally, Panel (c) tracks the perceptual result: the intensity of the second harmonic (H2) grows over time in response to this shift. This "bloom" is a defining characteristic of high-quality tube simulations, providing a dynamic spectral texture that evolves in direct response to the player's touch and sustain.

Figure 4 --- Signal Comparison and Blocking Distortion. Figure 4 provides a detailed time-domain analysis of the audible phenomena produced by the model. Panel (a) shows the response to a decaying note over 200 ms, where the "strangling" of positive peaks is clearly visible compared to a static $\tanh$ model. The detail views in panels (b) and (c) highlight the evolution of the distortion profile: the initial attack (b) shows the immediate attenuation of positive excursions (instantaneous conduction), while the sustained phase (c) reveals a significant asymmetric vertical shift due to the accumulated $\Delta b(t)$. Panel (d) plots the "deformation intensity" (the absolute difference between the static and dynamic models), illustrating how our model provides a time-varying texture that is fundamentally different from a memoryless waveshaper. The key observation is the asymmetric compression visible in the output: positive peaks (which triggered conduction) are attenuated and distorted differently from negative peaks, producing a visually and auditorily distinct asymmetric waveform consistent with measured tube amplifier behavior \cite{pakarinen2010distortion}.

Figure 5 — Harmonic Spectrum. The FFT comparison in Figure 5 provides a frequency-domain view of the model's effect. Without grid current (left panel), the static tanh model generates predominantly odd harmonics (H3, H5, H7...) due to its point-symmetric transfer function. With grid current active (right panel), the dynamic bias shift breaks the odd symmetry, enhancing even harmonics — particularly H2 and H4 — by up to 15–20 dB relative to the no-GC case. This behavior aligns precisely with measurements reported in [CITE] for real JCM800 and 5150 tube stages, validating the physical relevance of the model despite its simplified parametric form.

Figure 6 --- Dynamic Recovery and the "Breathe" Effect. Figure 6 demonstrates the recovery phase of the amplifier after a high-amplitude transient (e.g., a heavy pick attack or power chord), utilizing "Ultimate Contrast" visual markers for improved clarity. Panel (a) shows the output waveform for a loud burst followed by a sustained signal at a level (0.2) below the conduction threshold ($V_{th} = 0.25$). This isolates the recovery phase from further charge accumulation. The "Fast Recovery" ($\tau = 5$ ms) is shown as a dotted green line, returning instantly to symmetry. The "Slow Breathe" setting ($\tau = 800$ ms) is shown as a solid red line, remaining heavily choked and asymmetrical throughout the sustain. Panel (b) reveals the physical cause of this phenomenon: the DC bias shift (dotted black line) slowly "crawls" back toward zero, a recovery trajectory that perfectly correlates with the rebound of the auditory volume (RMS envelope). This visualization providing a clear, pedagogical link between the internal state of the virtual coupling capacitor and the resulting musical "breathing" of the preamp.

3.2.4 Model Validity and Limitations
The proposed model faithfully reproduces the following physically observed phenomena, within the constraints of real-time audio processing:

Threshold onset of grid conduction, approximated by half-wave rectification (Eq. 1), consistent with the tube's exponential $i_g$–$V_{gk}$ characteristic for $V_{gk} \gtrsim 0$.
RC charge–discharge dynamics of the coupling network (Eq. 2), implemented with negligible CPU overhead as a single-pole IIR filter.
Input attenuation (Eq. 3) producing amplitude-dependent "strangling" consistent with the known guitar amplifier phenomenon of blocking distortion.
Dynamic bias shift and harmonic enrichment (Eqs. 4–5), generating dynamically varying even-order harmonics in agreement with measured spectra.
The model has three main simplifications relative to a full circuit simulation. First, the rectification function $\max(0, x - V_{th})$ is a piecewise-linear approximation of the exponential $i_g$–$V_{gk}$ diode characteristic; this difference is perceptually negligible given the subsequent RC smoothing and the strong nonlinearity of the tanh stage. Second, the model is memoryless within each stage with respect to the bias shift path, whereas in a real circuit the bias perturbation would feed back into the grid voltage through the coupling network; this feedback loop is approximated by the temporal accumulation of the RC envelope. Third, the model does not account for screen grid or plate current saturation effects, which are partially captured by the global tanh stage but not individually parametrized.

Despite these simplifications, the model achieves $\mathcal{O}(1)$ complexity per sample, compilable to efficient WebAssembly via the FAUST compiler, and perceptually indistinguishable from reference recordings in informal listening tests. Its four parameters ($V_{th}$, $\tau$, $\alpha_{in}$, $\alpha_{bias}$) map directly to physically meaningful quantities of the tube circuit, facilitating preset design for specific amplifier models and expressive real-time control during performance.

